\documentclass[10pt,a4paper]{amsart}
\usepackage[margin=19mm]{geometry}
\usepackage{amsmath,amssymb,mathtools}
\usepackage[scr=rsfs,cal=euler]{mathalfa}
\usepackage{xfrac}
\usepackage[unicode,hidelinks,bookmarks=false]{hyperref}
\newcommand{\ie}{i.e.\ }
\newcommand{\cf}{cf.\ }
\newcommand{\resp}{respectively\ }
\newcommand{\Subsection}{\subsection}
\providecommand{\nobreakdash}{\nobreak\hbox{-}}
\setlength{\emergencystretch}{2em}
\multlinegap0pt
\usepackage{amssymb}
\newtheorem{thm}{Theorem}
\newcommand{\D}{{\mathbb D}}
\def\be{\begin{equation}}
\def\ee{\end{equation}}
\title{On some properties of logarithmic coefficients of inverse of univalent functions}
\author{Milutin Obradovi\'c, Nikola Tuneski, Pawe{\l} Zaprawa}
\date{Source: arXiv:2605.14425v1 (CC BY 4.0)}
\begin{document}
\maketitle

\noindent\emph{Source and licence.} On some properties of logarithmic coefficients of inverse of univalent functions. Version arXiv:2605.14425v1 (CC BY 4.0), distributed by arXiv under the Creative Commons Attribution 4.0 International licence, http://creativecommons.org/licenses/by/4.0/, as recorded in the arXiv licence metadata for this exact version and shown on the abstract page https://arxiv.org/abs/2605.14425v1. Full licence notice: \texttt{licenses/arxiv-2605.14425-CC-BY-4.0.txt}.\par\smallskip

\noindent\textit{Editorial scope.} Counted unit: the article's thm th1 (source0.tex lines 307--315). The article's complete original proof of it is delivered with it (source0.tex lines 317--349, one \texttt{proof} environment of 33 lines). No mathematics is written, repaired or re-derived: the statement and the proof are the article's own lines, in the article's own notation, and no retained line is substituted or re-flowed.  Retained uncounted as the source-local context the proof needs: lines 166--170; lines 183--189; lines 201--203; lines 205--211; lines 265--267; lines 271--278; lines 287--289; lines 291--293.  Nothing was omitted from any retained span, so the package carries no omission record.  No retained line contains a citation, so the wrapper carries no bibliography and no bibliography entry.  No retained line contains a figure environment, an image-inclusion command or drawing code, so no image is delivered and the package declares no asset dependency.  Editorial formatting only, in addition: an AMS \texttt{amsart} preamble that restates the article's own declarations and macros used by the retained lines, namely the environments \texttt{thm} on separate counters, together with the standard packages those lines need.  The retained environments print in this document as Theorem 1 in order of appearance, whereas the article prints the same items with the numbering of its own full document.  The complete native source of the article as distributed in the arXiv e-print for this exact version is bundled unchanged as \texttt{sources/200584/source0.tex}.
Let $\mathcal{A}$ be the class of functions $f$ analytic in the open unit disk $\D=\{z:|z|<1\}$ of the form
\be\label{eq-1}
f(z)=z+a_2z^2+a_3z^3+\cdots,
\ee
and let $\mathcal{S}$ be the subclass of $\mathcal{A}$ consisting of functions that are univalent in $\D$.

\begin{equation}\label{eq-3}
\begin{split}
 \gamma_{1}&=\frac{a_{2}}{2},\\
 \gamma_{2}&=\frac{1}{2}\left(a_3-\frac{1}{2}a_{2}^{2}\right), \\
 \gamma_{3}&=\frac{1}{2}\left(a_4-a_2a_3+\frac13a_2^3\right).
\end{split}
\end{equation}

\begin{equation}\label{eq-4}
f^{-1}(w) = w+A_2w^2+A_3w^3+\cdots.
\end{equation}

\begin{equation}\label{eq-5}
\begin{array}{l}
A_{2}=-a_{2}, \\
A _{3}=-a_{3}+2a_{2}^{2} , \\
A_{4}= -a_{4}+5a_{2}a_{3}-5a_{2}^{3}.
\end{array}
\end{equation}

\begin{equation}\label{eq-10a}
|\omega_{13}|^{2}+3|\omega_{33}|^{2}\leq\frac13.
\end{equation}

\begin{equation}\label{eq-11}
\begin{split}
a_{2}&=2\omega _{11},\\
a_{3}&=2\omega_{13}+3\omega_{11}^{2}, \\
a_{4}&=2\omega_{33}+8\omega_{11}\omega_{13}+\frac{10}{3}\omega_{11}^{3},\\
0&=3\omega_{15}-3\omega_{11}\omega_{13}+\omega_{11}^{3}-3\omega_{33}.
\end{split}
\end{equation}

\begin{equation}\label{eq-99}
|a_{3}-a_{2}^{2}|\leq1
\end{equation}

\begin{equation}\label{eq-100}
|2\omega_{13}-\omega_{11}^{2}|\leq1.
\end{equation}

\begin{thm}\label{th1}
Let $f\in\mathcal{S}$ and let $\Gamma_{1},\Gamma_{2},\Gamma_{3}$ be initial logarithmic coefficients of
its inverse function $f^{-1}$ given by \eqref{eq-5}. Then the following sharp inequalities hold
\begin{itemize}
  \item[($i$)] $|\Gamma_{1}|\leq 1$;
  \item[($ii$)] $|\Gamma_{2}| \leq\frac{3}{2}$;
  \item [($iii$)]$|\Gamma_{3}|\leq\frac{10}{3}$.
\end{itemize}
\end{thm}

\begin{proof}
Let $f\in\mathcal{S}$ is given by \eqref{eq-1} and its inverse $f^{-1}$ is given by \eqref{eq-4}, then using the relations \eqref{eq-3} and \eqref{eq-5}, we have
\[ |\Gamma_{1}|=\frac{1}{2}|A_{2}|=\frac{1}{2}|-a_{2}|\leq 1\]
and
\[|\Gamma_{2}|=\frac{1}{2}|A_{3}-\frac{1}{2}A_{2}^{2}|=\frac{1}{2}|-a_{3}+\frac{3}{2}a_{2}^{2}|
\leq \frac{1}{2}(|a_{3}-a_{2}^{2}|+\frac{1}{2}|a_{2}|^{2})\leq\frac{3}{2},\]
where we used \eqref{eq-99}.

\medskip

This proved (i) and (ii).

\medskip

For proving (iii), similarly, we obtain
\[
\Gamma_{3}=\frac{1}{2}\left( A_4-A_2A_3+\frac{1}{3}A_2^3\right)=\frac{1}{2}\left(-a_{4}+4a_{2}a_{3}-\frac{10}{3}a_{2}^{3}\right),
\]
or if we use the Grunsky coefficients given by relations \eqref{eq-11}:
\[\Gamma_{3} = -\omega_{33}+4\omega_{11}\omega_{13}-3\omega_{11}^{3} =-\omega_{33}+2\omega_{11}(2\omega_{13}-\omega_{11}^2)-\omega_{11}^{3}.\]
From \eqref{eq-100} and \eqref{eq-10a}, we have
\[
\begin{split}
|\Gamma_{3}| &\leq |\omega_{33}|+2|\omega_{11}|+|\omega_{11}|^{3} \\
&\leq \frac13\sqrt{1-3|\omega_{13}|^2}+2|\omega_{11}|+|\omega_{11}|^{3} \\
&\leq \frac{10}{3}\ .
\end{split}
\]

\medskip

All three estimates are sharp as the Koebe function $k(z)=\frac{z}{(1-z)^{2}}$ shows.
\end{proof}

\end{document}
